\FloatBarrier
\section{Conclusiones}\label{sec:conclusiones}

\subsection{Conclusiones del equipo}\label{sec:conclusiones-equipo}

% Propuesta de conclusión general; pendiente de revisión y acuerdo del equipo.
En conjunto, el desarrollo del Sistema de Control de Laboratorio de Calidad de Harina de Trigo nos permitió aplicar los conocimientos de ingeniería de software a una necesidad concreta: organizar la información de clientes, equipos, lotes e inspecciones para facilitar el seguimiento de los resultados y la emisión de certificados. La división del trabajo por módulos, el uso de un repositorio compartido y las revisiones continuas nos ayudaron a entender la importancia de la comunicación, la integración y las pruebas durante todo el proceso. También identificamos que un sistema no se considera terminado solo porque sus pantallas funcionen de manera individual; es necesario comprobar que sus componentes intercambien información de forma correcta, que se respeten los permisos y que los datos puedan consultarse con confianza. Como equipo, este proyecto nos dejó una visión más completa del desarrollo de software, desde el análisis de requisitos y el diseño hasta la implementación, documentación y validación.

Desde la perspectiva de la generación de reportes y certificados, el proyecto demuestra que centralizar el registro de lotes, inspecciones y límites por cliente o norma permite emitir resultados más consistentes, trazables y fáciles de validar. Contar con procesos claros y documentación estructurada reduce ambigüedades al consolidar la información y facilita que los reportes cumplan su función real de apoyo a la toma de decisiones y al control de calidad.

\subsection{Lecciones aprendidas}\label{sec:lecciones}

Cada integrante escribe una lección propia.

\pend{una lección por integrante, escrita por cada persona, sin nombres en el repositorio}

% Propuesta de lección individual de Nathalia; pendiente de su revisión.
Al trabajar en la integración de sesiones, permisos y módulos operativos como Equipos, Lotes e Inspecciones, comprendí la importancia de diseñar componentes que no solo funcionen por separado, sino que respeten las mismas reglas de negocio y mantengan la información consistente. La coordinación entre módulos, el control de acceso y las validaciones me ayudaron a reconocer que muchos de los retos aparecen cuando se conectan procesos distintos. Esta experiencia fortaleció mi comprensión del desarrollo colaborativo y de la necesidad de probar cada cambio antes de incorporarlo al sistema completo.

% Lección individual correspondiente a Reportes (aportada previamente).
Al encargarme de la parte de reportes, comprendí que entender a fondo la documentación y el contexto del sistema es fundamental antes de construir o consolidar cualquier salida. Los procesos y flujos definidos en el proyecto existen precisamente para facilitar este trabajo y hacer que el funcionamiento del sistema sea claro y comprensible; cuando la documentación y los procesos se asimilan bien, los reportes resultan mucho más coherentes, confiables y fáciles de interpretar.
